\chapter{Zespolone funkcje elementarne}
\label{ma2:ch:complex-elementary}

Wykorzystujemy konstrukcję funkcji wykładniczej i trygonometrycznych z AM I,
rozdział 18, oraz twierdzenie o różniczkowaniu szeregów potęgowych
z rozdziału \ref{ma2:ch:power-series}. W szczególności odróżniamy funkcję
jednowartościową od zbioru wszystkich wartości logarytmu lub potęgi.

\section{Funkcja wykładnicza}

Przypomnijmy definicję AM I, 18.1.1:
\[
 \exp z=\sum_{n=0}^\infty\frac{z^n}{n!},\qquad z\in\C.
\]
Zgodnie z twierdzeniem AM I, 18.1.2, funkcja ta jest ciągła,
$\exp0=1$, $\exp(z+w)=\exp z\exp w$ i $\exp z\ne0$ dla wszystkich
$z,w\in\C$. Dowód znajduje się w tomie I. Ograniczenie do $\R$
jest rosnącą bijekcją na $(0,\infty)$ z ciągłą odwrotnością $\log$
(AM I, 18.1.3). Wzór Eulera dla rzeczywistego $t$,
$\exp(it)=\cos t+i\sin t$, wynika z rozdzielenia szeregu na wyrazy
parzyste i nieparzyste (AM I, 18.3).

\begin{twierdzenie}
\label{ma2:thm:complex-exp}
Funkcja $\exp$ jest różniczkowalna w sensie zespolonym na $\C$,
$\exp' z=\exp z$. Dla $x,y\in\R$ zachodzi
\[
 \exp(x+iy)=e^x(\cos y+i\sin y),\qquad |\exp(x+iy)|=e^x.
\]
Ponadto $\exp(\C)=\C\setminus\{0\}$, a
\[
 \exp z=\exp w\quad\Longleftrightarrow\quad z-w\in2\pi i\Z.
\]
\end{twierdzenie}
\begin{proof}
Iloraz kolejnych dodatnich współczynników szeregu wynosi $1/(n+1)$,
więc promień zbieżności jest nieskończony. Twierdzenie
\ref{ma2:thm:power-derivative} daje pochodną
$\sum_{n=1}^\infty nz^{n-1}/n!=\exp z$.
Prawo dodawania i wzór Eulera dają pierwszą równość, a
$\cos^2y+\sin^2y=1$ (AM I, 18.3.2) drugą.
Dla $v\ne0$ postać trygonometryczna z AM I, 18.8.1, daje
$v=r(\cos t+i\sin t)$ dla $r=|v|>0$ i $t\in\R$.
Wtedy $\exp(\log r+it)=v$, co dowodzi surjektywności na wskazany zbiór.
Jeśli $\exp z=\exp w$, to $\exp(z-w)=1$. Dla $z-w=x+iy$
równość modułów daje $e^x=1$, a zatem $x=0$. Jednoznaczność
argumentu modulo $2\pi$ (AM I, 18.8.1) daje $y\in2\pi\Z$.
Odwrotnie, dla $y\in2\pi\Z$ wzór Eulera daje $\exp(iy)=1$,
więc prawo dodawania daje $\exp(w+iy)=\exp w$.
\end{proof}

\section{Funkcje trygonometryczne i hiperboliczne}

\begin{definicja}
\label{ma2:def:complex-trig}
Dla $z\in\C$ definiujemy
\[
 \cos z=\frac{e^{iz}+e^{-iz}}2,\quad
 \sin z=\frac{e^{iz}-e^{-iz}}{2i},\qquad
 \cosh z=\frac{e^z+e^{-z}}2,\quad
 \sinh z=\frac{e^z-e^{-z}}2 .
\]
Na zbiorach, gdzie odpowiednie mianowniki są niezerowe, określamy
$\tan z=\sin z/\cos z$, $\cot z=\cos z/\sin z$,
$\tanh z=\sinh z/\cosh z$ i $\coth z=\cosh z/\sinh z$.
\end{definicja}

\begin{twierdzenie}
\label{ma2:thm:complex-trig-identities}
Funkcje $\cos,\sin,\cosh,\sinh$ są różniczkowalne w sensie
zespolonym na $\C$. Ich szeregi mają nieskończony promień zbieżności:
\begin{align*}
 \cos z&=\sum_{n=0}^\infty\frac{(-1)^nz^{2n}}{(2n)!},
 &\sin z&=\sum_{n=0}^\infty\frac{(-1)^nz^{2n+1}}{(2n+1)!},\\
 \cosh z&=\sum_{n=0}^\infty\frac{z^{2n}}{(2n)!},
 &\sinh z&=\sum_{n=0}^\infty\frac{z^{2n+1}}{(2n+1)!}.
\end{align*}
Mamy $\sin'=\cos$, $\cos'=-\sin$, $\sinh'=\cosh$, $\cosh'=\sinh$,
a dla $z,w\in\C$:
\begin{align*}
 \sin(z+w)&=\sin z\cos w+\cos z\sin w,\\
 \cos(z+w)&=\cos z\cos w-\sin z\sin w,\\
 \sinh(z+w)&=\sinh z\cosh w+\cosh z\sinh w,\\
 \cosh(z+w)&=\cosh z\cosh w+\sinh z\sinh w.
\end{align*}
W szczególności $\cos^2z+\sin^2z=1$ i $\cosh^2z-\sinh^2z=1$.
\end{twierdzenie}
\begin{proof}
Podstawiamy szeregi wykładnicze do definicji; bezwzględna
zbieżność pozwala zebrać wyrazy parzyste i nieparzyste. Sumy
szeregów modułów są ograniczone przez $\exp|z|$. Promienie są
więc nieskończone. Różniczkowanie wyraz po wyrazie na podstawie
twierdzenia \ref{ma2:thm:power-derivative} i przesunięcie indeksu
dają podane cztery pochodne. Dla wzorów dodawania piszemy
$A=e^{iz}$, $B=e^{iw}$, a następnie
\[
\begin{split}
 \sin z\cos w+\cos z\sin w
 &=\frac{(A-A^{-1})(B+B^{-1})+(A+A^{-1})(B-B^{-1})}{4i}\\
 &=\frac{AB-(AB)^{-1}}{2i}=\sin(z+w),\\
 \cos z\cos w-\sin z\sin w
 &=\frac{(A+A^{-1})(B+B^{-1})+(A-A^{-1})(B-B^{-1})}{4}\\
 &=\frac{AB+(AB)^{-1}}2=\cos(z+w).
\end{split}
\]
Dla hiperbolicznych bierzemy $A=e^z$, $B=e^w$: suma dwóch
iloczynów odpowiednio z jedną różnicą daje
$(AB-(AB)^{-1})/2$, a suma iloczynów sum i różnic daje
$(AB+(AB)^{-1})/2$. Wynik jest konsekwencją prawa dodawania
funkcji wykładniczej. Podstawienie $w=-z$ oraz parzystość
cosinusów i nieparzystość sinusów (z szeregów) dają tożsamości
kwadratów. Na $\R$ otrzymujemy funkcje z AM I, 18.3 i 18.6.
\end{proof}

\begin{twierdzenie}
\label{ma2:thm:complex-trig-zeros}
Zbiory zer funkcji $\sin,\cos,\sinh,\cosh$ wynoszą odpowiednio
\[
 \pi\Z,\quad \pi/2+\pi\Z,\quad i\pi\Z,\quad i\pi/2+i\pi\Z.
\]
Na swoich dziedzinach funkcje ilorazowe są różniczkowalne
w sensie zespolonym i
\[
 \tan' z=\frac1{\cos^2z},\quad
 \cot' z=-\frac1{\sin^2z},\quad
 \tanh' z=\frac1{\cosh^2z},\quad
 \coth' z=-\frac1{\sinh^2z}.
\]
\end{twierdzenie}
\begin{proof}
Z definicji i niezerowości wykładniczej $\sin z=0$ wtedy i tylko
wtedy, gdy $e^{2iz}=1$, czyli $2iz\in2\pi i\Z$ na mocy
twierdzenia \ref{ma2:thm:complex-exp}. Analogicznie $\cos z=0$
wtedy i tylko wtedy, gdy $e^{2iz}=-1=e^{i\pi}$, czyli
$2iz-i\pi\in2\pi i\Z$. Dla funkcji hiperbolicznych odpowiednie
równania mają postać $e^{2z}=1$ i $e^{2z}=e^{i\pi}$;
rozwiązujemy je tym samym kryterium.
Wzór na pochodną ilorazu wynika z rozkładu ilorazu różnicowego:
dla $g(z)\ne0$ i małego $h$, $g(z+h)\ne0$ dzięki ciągłości,
a
\[
 \frac{(f/g)(z+h)-(f/g)(z)}{h}
 =\frac{g(z)\frac{f(z+h)-f(z)}h-f(z)\frac{g(z+h)-g(z)}h}
 {g(z+h)g(z)}.
\]
Przechodzimy do granicy i podstawiamy pochodne oraz tożsamości
kwadratów z twierdzenia \ref{ma2:thm:complex-trig-identities}.
\end{proof}

\begin{przyklad}
Dla rzeczywistych $x,y$, wzory dodawania oraz definicje dają
$\sin(x+iy)=\sin x\cosh y+i\cos x\sinh y$ i
$\cos(x+iy)=\cos x\cosh y-i\sin x\sinh y$.
W szczególności $|\sin(iy)|=|\sinh y|$ i $|\cos(iy)|=\cosh y$,
więc żadna z tych funkcji nie jest ograniczona na $\C$.
Nie przeczy to oszacowaniom na osi rzeczywistej.
\end{przyklad}

\section{Argument i logarytm zespolony}

Dla $z\ne0$ wszystkie rozwiązania równania $e^w=z$ tworzą
zbiór
\[
 \{\log|z|+i(t+2k\pi):k\in\Z\},
\]
gdzie $t$ jest dowolnym argumentem $z$. Wynika to z twierdzenia
\ref{ma2:thm:complex-exp}. Zapis zbioru nie oznacza funkcji
jednowartościowej.

\begin{definicja}
\label{ma2:def:log-branch}
Gałęzią logarytmu na zbiorze otwartym $U\subset\C\setminus\{0\}$
nazywamy ciągłą funkcję $L:U\to\C$ spełniającą $e^{L(z)}=z$.
\end{definicja}

\begin{twierdzenie}
\label{ma2:thm:principal-log}
Na zbiorze $D=\C\setminus(-\infty,0]$ istnieje dokładnie jedna
gałąź logarytmu o częściach urojonych w $(-\pi,\pi)$.
Ma postać
\[
 \operatorname{Log}z=\log|z|+i\operatorname{Arg}z,\qquad
 \operatorname{Arg}(x+iy)=2\arctan\frac{y}{|z|+x}.
\]
Funkcja $\operatorname{Log}$ jest ciągła, a na dodatniej osi
rzeczywistej zgadza się z rzeczywistym logarytmem.
\end{twierdzenie}
\begin{proof}
Dla $z=x+iy\in D$ mamy $r=|z|>0$ i $r+x>0$: równość
$r+x=0$ wymagałaby $y=0$ i $x\le0$. Niech
$u=y/(r+x)$ oraz $t=2\arctan u\in(-\pi,\pi)$.
Wzory podwojenia i $\tan(t/2)=u$ dają
$\cos t=(1-u^2)/(1+u^2)$, $\sin t=2u/(1+u^2)$.
Ponieważ $y^2=r^2-x^2$, mamy
\[
 1+u^2=\frac{2r}{r+x},\qquad 1-u^2=\frac{2x}{r+x}.
\]
Stąd $\cos t=x/r$ i $\sin t=y/r$. Zatem
$e^{\log r+it}=z$. Ciągłość wynika z ciągłości modułu,
logarytmu rzeczywistego i arcus tangensa (AM I, 18.1.3,
18.5.2), oraz niezerowości mianownika. Dwie wartości
logarytmu różnią się o $2k\pi i$; gdy obie części urojone są
w $(-\pi,\pi)$, ich różnica ma moduł mniejszy niż $2\pi$,
więc $k=0$. Dla $z>0$ argument z podanego wzoru wynosi $0$.
\end{proof}

\begin{twierdzenie}
\label{ma2:thm:log-derivative}
Każda gałąź logarytmu $L$ na zbiorze otwartym
$U\subset\C\setminus\{0\}$ jest różniczkowalna w sensie
zespolonym i $L'(z)=1/z$.
\end{twierdzenie}
\begin{proof}
Ustalmy $z\in U$ i niech $h\ne0$ dąży do zera tak, aby $z+h\in U$.
Oznaczmy $\eta=L(z+h)-L(z)$. Z ciągłości $L$ mamy $\eta\to0$.
Ponadto $e^\eta=(z+h)/z$, więc $\eta\ne0$ dla $h\ne0$.
Stąd
\[
 \frac{L(z+h)-L(z)}h
 =\frac1z\,\frac{\eta}{e^\eta-1}\longrightarrow\frac1z,
\]
ponieważ pochodna wykładniczej w zerze wynosi $1$.
Nie zakładaliśmy wcześniej różniczkowalności funkcji odwrotnej.
\end{proof}

\begin{twierdzenie}
\label{ma2:thm:log-branches}
Na spójnym zbiorze otwartym dwie gałęzie logarytmu różnią się
o stałą $2k\pi i$, $k\in\Z$. Dla każdego $\theta\in\R$
na zbiorze $D_\theta=e^{i\theta}D$ istnieje gałąź
$L_\theta(z)=i\theta+\operatorname{Log}(e^{-i\theta}z)$.
Nie istnieje ciągła gałąź logarytmu na całym
$\C\setminus\{0\}$.
\end{twierdzenie}
\begin{proof}
Różnica gałęzi jest ciągła i przyjmuje wartości w dyskretnym
zbiorze $2\pi i\Z$ zgodnie z twierdzeniem
\ref{ma2:thm:complex-exp}. Obraz zbioru spójnego jest spójny;
jedynymi niepustymi spójnymi podzbiorami tego zbioru dyskretnego
są singletony. Wzór na $L_\theta$ jest ciągły i jego
wykładnicza wynosi $e^{i\theta}e^{-i\theta}z=z$.
Gdyby istniała globalna gałąź $L$, funkcja
$t\mapsto L(e^{it})-it$ na przedziale $[0,2\pi]$ byłaby ciągła
i przyjmowałaby wartości w $2\pi i\Z$, a więc byłaby stała.
Jednak jej wartości na końcach różnią się o $-2\pi i$,
co jest sprzecznością.
\end{proof}

\begin{przyklad}
Dla $z=w=e^{3\pi i/4}$ wszystkie trzy liczby $z,w,zw$
należą do $D$, lecz $\operatorname{Log}(zw)=-i\pi/2$,
natomiast $\operatorname{Log}z+\operatorname{Log}w=3\pi i/2$.
Ogólnie dla $z,w,zw\in D$ różnica tych wyrażeń należy do
$2\pi i\Z$, gdyż ich wykładnicze są równe. Dokładna równość
zachodzi wtedy i tylko wtedy, gdy
$\operatorname{Arg}z+\operatorname{Arg}w\in(-\pi,\pi)$:
wtedy suma jest logarytmem o wymaganej części urojonej,
a odwrotnie część urojona $\operatorname{Log}(zw)$ leży
w tym przedziale.
\end{przyklad}

\section{Potęgi i pierwiastki}

\begin{definicja}
\label{ma2:def:complex-power}
Dla $\alpha\in\C$ i gałęzi $L$ na $U$ określamy
$P_{\alpha,L}(z)=\exp(\alpha L(z))$, $z\in U$.
W przypadku $L=\operatorname{Log}$ mówimy o wartości głównej
potęgi na $D$. Wszystkie wartości potęgi w punkcie $z\ne0$
otrzymujemy jako $\exp(\alpha(w+2k\pi i))$, $k\in\Z$,
gdzie $e^w=z$.
\end{definicja}

\begin{twierdzenie}
\label{ma2:thm:complex-power}
Dla ustalonej gałęzi $L$ i $\alpha,\beta\in\C$ zachodzi
\[
 P_{\alpha,L}P_{\beta,L}=P_{\alpha+\beta,L},\qquad
 P'_{\alpha,L}=\alpha P_{\alpha-1,L}.
\]
Dla $m\in\Z$ funkcja $P_{m,L}$ zgadza się ze zwykłą
potęgą całkowitą i nie zależy od wyboru gałęzi.
Zmiana $L$ na $L+2k\pi i$ mnoży $P_{\alpha,L}$ przez
$e^{2k\pi i\alpha}$. Wszystkie wartości potęgi w danym
punkcie są jednakowe wtedy i tylko wtedy, gdy $\alpha\in\Z$.
\end{twierdzenie}
\begin{proof}
Pierwsza równość wynika z prawa dodawania wykładniczej.
Dla pochodnej stosujemy regułę łańcuchową do wykładniczej
i $L$ (twierdzenia \ref{ma2:thm:complex-exp}
i \ref{ma2:thm:log-derivative}), otrzymując
$P'_{\alpha,L}(z)=\alpha e^{\alpha L(z)}/z
=\alpha e^{(\alpha-1)L(z)}$. Reguła łańcuchowa w sensie
zespolonym wynika z rozwinięć
$g(z+h)=g(z)+g'(z)h+o(|h|)$ i
$f(w+v)=f(w)+f'(w)v+o(|v|)$:
dla $w=g(z)$, $v=g(z+h)-g(z)=O(|h|)$
reszta z drugiego rozwinięcia jest $o(|h|)$ (również
gdy $v=0$), a współczynnik przy $h$ wynosi $f'(w)g'(z)$.
Dla $m\ge0$ stosujemy prawo wykładnicze indukcyjnie,
a dla $m<0$ korzystamy z $e^{-L(z)}=1/z$.
Wzór zmiany gałęzi znów wynika z prawa dodawania.
Jeżeli wszystkie wartości są jednakowe, porównanie
$k=0$ i $k=1$ daje $e^{2\pi i\alpha}=1$, gdyż
wykładnicza nie znika. Z twierdzenia
\ref{ma2:thm:complex-exp} wynika $\alpha\in\Z$.
Implikacja odwrotna została już udowodniona.
\end{proof}

\begin{przyklad}
Dla $n\ge1$ całkowitego, $P_{1/n,\operatorname{Log}}$
jest ciągłym wyborem pierwiastka na $D$. Istotnie,
jego $n$-ta potęga wynosi $z$. Wszystkie pierwiastki
niezerowego $z$ są postaci
\[
 P_{1/n,\operatorname{Log}}(z)e^{2k\pi i/n},\quad
 k=0,\ldots,n-1,\qquad z\in D,
\]
zgodnie z AM I, 18.8.2. Są różne na mocy kryterium
równości wykładniczych. Dla $z=0$ jedynym pierwiastkiem
jest $0$. Natomiast dla $\alpha=i$ i $z=1$ wszystkie
wartości to $e^{-2k\pi}$, $k\in\Z$; jest ich nieskończenie wiele.
\end{przyklad}

\begin{uwaga}
Tożsamości potęg przy różnych gałęziach, w szczególności
$(z^\alpha)^\beta=z^{\alpha\beta}$, nie są automatyczne.
Na przykład dla $z=e^{3\pi i/4}$ i wartości głównych,
$(z^2)^{1/2}=e^{-i\pi/4}\ne z$. W tym przykładzie zarówno
$z$, jak i $z^2=-i$ należą do $D$. Funkcje odwrotne
do zespolonego sinusa i cosinusa wymagają dodatkowych wyborów
gałęzi; ich globalna teoria nie jest potrzebna w dalszym
rozdziale o funkcjach analitycznych.
\end{uwaga}
